\section{A shared-KV store}
\label{sec:sharedkv}
\label{sec:csa2}

\secref{sec:longcontext} makes retrieval $O(1)$ in \emph{slots}---a query needs
one or two---but the scan that \emph{finds} those slots is still $O(M)$: the read
scores $q\!\cdot\!K$ over all $M$ keys. DeepSeek-V4.1-Flash (CSA2) makes
long-context key--value cheap by exploiting inter-layer redundancy along three
axes: a sparse candidate pool computed once and reused across layers, cross-layer
KV reuse (a few anchor layers compute, the rest reuse), and low-bit (FP4) KV. We
port all three to the store as an eval-time read path---no retraining---and each
preserves recall.

\subsection{Candidate pool: $O(1)$-slot becomes $O(1)$-compute}

We restrict the read to an index-selected pool of candidate slots rather than the
whole store. Two indices:
\begin{itemize}
\item \textbf{token} (exact). We record each slot's key-token at write; a query
retrieves the slots whose stored token equals its own. For the position-free
layer-0 key this is an \emph{exact} hash.
\item \textbf{ann} (approximate). Random-hyperplane LSH on the key vectors buckets
slots; a query reads its bucket, which works when keys are fuzzy rather than exact
tokens.
\end{itemize}
\tabref{tab:pool} shows the token index scores $\approx\!2$ of $8001$ slots at
$63$ windows with recall equal to the dense scan; the LSH index scores $\approx
M/12$, also equal to dense. Restricting the read to the pool makes it $O(\text{pool})$,
independent of $M$.

\begin{table}[t]
\centering
\begin{tabular}{lccc}
\toprule
store (63 windows, 8001 slots) & dense (\texttt{none}) & token & ann (LSH) \\
\midrule
candidates scored & 8001 & \textbf{1.9} & 642 \\
recall & \recall{82.5} & \recall{81.9} & \recall{83.1} \\
\bottomrule
\end{tabular}
\caption{Candidate-pool read: the token index scores $\approx2$ of $8001$ slots at
dense-matching recall---the $O(1)$-slot read becomes $O(1)$-compute---and LSH gives
a $\sim\!10\times$ cut for fuzzy keys. \texttt{none} reproduces the dense read of
\S\ref{sec:longcontext}.}
\label{tab:pool}
\end{table}

\subsection{Cross-layer reuse is free}

When the model reads at several layers, the pool need not be rebuilt per layer.
Building it once at an anchor layer and reusing the same candidates at the later
read layers gives \emph{identical} recall to rebuilding per layer (\recall{100}
either way on a six-read-layer model), at one pool build instead of six---exactly
CSA2's ``few anchor layers, the rest reuse.'' This works because, as
\S\ref{sec:placement} shows independently, the layers are redundant enough that one
shared retrieval serves them all.

\subsection{Low-bit KV}

Finally we quantize the stored codes. The store survives \textbf{4 bits per element
with no recall loss} ($8\times$ smaller than the fp32 store and $4\times$ smaller than fp16: $1024\!\to\!256$ bytes/slot against fp16; at
$63$ windows $4$-bit equals $16$-bit at \recall{82.5}), and $3$-bit is a hard
cliff. Ablating keys versus values, the \emph{keys} are the limit: matching needs
$4$ bits to keep near-tied keys apart, while values decode correctly at $3$ bits,
so an asymmetric keys-$4$/values-$3$ store packs tighter still.

\paragraph{Summary.} Together, the store is a shared key--value memory that holds
$\gg$ the context window ($63\times$), is retrieved in $O(1)$ slots at
$O(1)$-compute, is reused across layers for free, and costs $4$ bits per
element---the DeepSeek shared-KV construction realised on an in-stack episodic
store. We emphasise these are \emph{measured} properties: the probe still computes
$q\!\cdot\!K$ densely and fake-quantizes in fp32 to obtain the recall numbers;
turning them into literal speed and bytes needs a true hashed/int4 kernel, which is
engineering, not a further result.
